\section{Equilibration with an explicit residual budget}
\label{sec:contract}

Block ownership is not the only information a generator can supply. An application
can also state how much violation of each original constraint is acceptable. Such a
budget must come from the application or be declared as a numerical test parameter;
it cannot be inferred reliably from a zero right-hand side or a block label. Scaling
to make solver tolerances meaningful is established practice
\cite{gurobiScaling,idaesScaling}. The contribution of the following construction is
an explicit feasibility test and an optimality characterization for combining that
budget with coefficient limits, including a reason to abstain.

For a nonconstant row $a^Tx\le b$, define the original residual
$v(x)=(a^Tx-b)_+$ and specify a budget $\delta>0$. Let $\epsilon>0$ be a
\emph{hypothetical bound on the residual of the exported row}. A solver's feasibility
parameter alone does not establish this hypothesis: internal scaling, presolve,
integrality tolerances and finite precision intervene. For a positive row factor
$\alpha$, the identity
\begin{equation}\label{eq:residual-transfer}
 (\alpha a^Tx-\alpha b)_+=\alpha v(x)
\end{equation}
holds in exact arithmetic. Consequently, the hypothesis of exported residual at
most $\epsilon$ implies $v(x)\le\delta$ whenever $\alpha\ge\epsilon/\delta$.
Absolute residuals give the same identity for equality constraints.

Write $m=\min_{a_j\ne0}|a_j|$ and $M=\max_j|a_j|$, and specify coefficient
limits $0<\ell\le U$. These bounds concern the constraint matrix, not variable
domains, objective coefficients or the right-hand side. If an RHS magnitude limit
$U_b$ is also required, intersect the upper bound below with $U_b/|b|$ for $b\ne0$.

\begin{proposition}[Feasibility and optimal constrained centering]
\label{prop:budget}
The factors that place every nonzero coefficient in $[\ell,U]$ and satisfy
$\epsilon/\alpha\le\delta$ form exactly the interval
\begin{equation}\label{eq:budget-interval}
 \mathcal I=[L,H],\qquad
 L=\max\{\epsilon/\delta,\ell/m\},\qquad H=U/M.
\end{equation}
The interval is nonempty if and only if
\begin{equation}\label{eq:budget-feasibility}
 M/m\le U/\ell\quad\text{and}\quad \delta\ge\epsilon M/U.
\end{equation}
When it is nonempty, the unique minimizer of
\begin{equation}\label{eq:budget-objective}
 \min_{\alpha\in\mathcal I}\;\max_{a_j\ne0}|\log(\alpha|a_j|)|
\end{equation}
is
\begin{equation}\label{eq:budget-optimum}
 \alpha^*=\min\{H,\max\{L,(mM)^{-1/2}\}\}.
\end{equation}
If $L>H$, no positive row factor meets both requirements.
\end{proposition}

\begin{proof}
The coefficient bounds are equivalent to $\alpha m\ge\ell$ and $\alpha M\le U$.
Together with $\alpha\delta\ge\epsilon$ they give \eqref{eq:budget-interval}.
The condition $L\le H$ is equivalent to the two inequalities in
\eqref{eq:budget-feasibility}. To solve the objective, put $t=\log\alpha$,
$q=(\log m+\log M)/2$ and $h=(\log M-\log m)/2$. Since the extreme absolute
logarithms occur at an endpoint,
\[
 \max_j|t+\log|a_j||=\max\{|t+\log m|,|t+\log M|\}=h+|t+q|.
\]
The unique minimizer on $[\log L,\log H]$ is the projection of $-q$ onto that
interval. Exponentiating gives \eqref{eq:budget-optimum}.
\end{proof}

Thus the rule optimizes a coefficient-centering criterion; it does not minimize a
condition number or predict running time. The two causes of abstention have different
interpretations. A within-row ratio exceeding $U/\ell$ cannot be repaired by row
scaling. A budget below $\epsilon M/U$ conflicts with the prescribed coefficient
ceiling even if the row is perfectly balanced. Returning the base row does not
resolve either conflict and must not be reported as satisfying the contract.
The test identifies an incompatibility among numerical specifications, not an
infeasible optimization model. Revising the units, limits or formulation requires a
separate decision.

\begin{corollary}[Invariance under positive re-emission]
\label{cor:budget-invariance}
Replace $(a,b,\delta)$ by $(qa,qb,q\delta)$ for any $q>0$, keeping
$\epsilon,\ell,U$ fixed. The interval becomes $\mathcal I/q$, abstention is
unchanged, and the selected factor becomes $\alpha^*/q$. The exported row and the
budget-normalized original residual $v(x)/\delta$ are therefore unchanged.
\end{corollary}
\begin{proof}
Both coefficient extremes and the original budget acquire the factor $q$.
Each term defining $L$, $H$ and $(mM)^{-1/2}$ is divided by $q$.
Equation~\eqref{eq:residual-transfer} gives the residual statement.
\end{proof}

This invariance is lost if the row is re-emitted but its budget is not transformed
with it. It also distinguishes a budget in canonical units from a fixed cutoff on
the scaling factor itself. Restricting factors to powers of two can avoid additional
rounding in multiplication when no underflow or overflow occurs, but this familiar
technique \cite{berthold2021scaling} preserves the stated invariance only for
power-of-two re-emissions, subject to representability. The dyadic construction below
explicitly checks those arithmetic restrictions.

\subsection{The part that row scaling cannot fix: integer rounding}
Let $J$ be the integer columns and let $\bar x$ be obtained by rounding only these
coordinates of $x$. Suppose $|\bar x_j-x_j|\le\eta_j$ for $j\in J$ and put
$E=\sum_{j\in J}|a_j|\eta_j$. The triangle inequality and
\eqref{eq:residual-transfer} give the useful but conditional bound
\begin{equation}\label{eq:rounding-budget}
 v(\bar x)\le \epsilon/\alpha+E.
\end{equation}
It applies to equalities with absolute residuals as well. The two terms have
different behavior: changing the row factor can reduce the exported-residual term,
whereas it leaves the amplification of integer rounding in original units unchanged.
The scalar normalization of a big-$M$ row therefore cannot, by itself, remove
the effect of an integrality tolerance on that row.

For a contract based on the worst-case bound \eqref{eq:rounding-budget},
Proposition~\ref{prop:budget} applies after replacing $\delta$ by $\delta-E$,
provided $\delta>E$. If $\delta\le E$ and $\epsilon>0$, no positive factor
satisfies this worst-case contract. This does not say that every returned solution
will violate its budget: actual rounding errors may be smaller or cancel, and other
constraints can exclude the worst case. It says that the stated information is
insufficient to certify it. Re-emission invariance still holds because $E$ transforms
to $qE$. A tighter integrality bound or a reformulation can reduce $E$; pure row
scaling cannot. This distinction prevents the row-only guarantee from being
mistaken for a guarantee about a rounded MIP solution.

\subsection{An exact export on a dyadic grid}
Even when a factor satisfies the real-arithmetic interval, independently rounded
products $\operatorname{fl}(\alpha a_j)$ and $\operatorname{fl}(\alpha b)$ need not
describe an exactly proportional row over the stored data. Merely rounding a
computed endpoint to nearest can also move it outside a closed admissible interval.
Our arbitrary-factor implementation forms the interval endpoints as exact rational
numbers in the supplied binary64 inputs and rounds them \emph{inward}; it abstains
if the interval contains no positive binary64 factor. This certifies factor
admissibility, not exactness of all subsequent products.

Powers of two provide an established way to avoid multiplication rounding. The
following explicit intersection also accounts for subnormal numbers and for the
RHS. Assume finite binary64 input data and IEEE arithmetic with gradual underflow
\cite{ieee2019float}. Write each nonzero
datum $z\in\{a_1,\ldots,a_n,b\}$ uniquely as
$z=\operatorname{sign}(z)N_z2^{p_z}$ with $N_z$ a positive odd integer, and put
$q_z=\lfloor\log_2|z|\rfloor$. Let $[L,H]$ be the desired real factor interval,
including any rounding allowance from \eqref{eq:rounding-budget}.

\begin{proposition}[Exact dyadic admissibility]
\label{prop:dyadic}
A factor $\alpha=2^k$ is positive and representable in binary64, belongs to
$[L,H]$, and multiplies every nonzero coefficient and the RHS exactly into a
finite binary64 value if and only if the integer $k$ satisfies
\begin{align*}
 K_-&=\max\left\{-1074,\lceil\log_2L\rceil,
                    \max_z(-1074-p_z)\right\}\le k,\\
 k&\le K_+=\min\left\{1023,\lfloor\log_2H\rfloor,
                    \min_z(1023-q_z)\right\}.
\end{align*}
If $K_->K_+$ there is no such factor. Otherwise an exponent nearest to
$-\tfrac12\log_2(mM)$ among the integers in $[K_-,K_+]$ minimizes the
logarithmic coefficient objective \eqref{eq:budget-objective} on this exact-export
set. A tie may be resolved by choosing the larger factor.
\end{proposition}
\begin{proof}
The factor itself requires $-1074\le k\le1023$. Since the datum's odd
significand is unchanged, the product $N_z2^{p_z+k}$ lies on the subnormal grid
exactly when $p_z+k\ge-1074$; its highest bit stays below overflow exactly when
$q_z+k\le1023$. The original significand has at most 53 bits, so these two
conditions are also sufficient. Intersecting them over the row data and with
$L\le2^k\le H$ gives the interval of integer exponents. The objective is
$\tfrac12\log(M/m)+|k\log2+\tfrac12\log(mM)|$, proving the nearest-exponent
claim.
\end{proof}

The implementation obtains the endpoint exponents using integer ratios and
bit lengths, and compares candidate centering objectives exactly as rational numbers. It also
checks the rounded products against exact rational products.
A round-trip decimal serializer is necessary if the row is subsequently written
to a text file: fixed decimal places can discard small values even when the
multiplication is exact. These checks concern the represented mathematical model.
They still do not guarantee that the solver will return an acceptable point.

\paragraph{Verification after solving.}
The interval theorem is conditional and exact-arithmetic. Every returned point must
still be checked against the original rows, bounds and integrality restrictions.
Our small stress experiments compute row residuals by exact rational arithmetic on
the supplied binary64 coefficients and returned solution. They also enumerate the
binary assignments to determine the optimum independently. These checks are possible
because those instances are deliberately small. On large dispatch models the archived
checker uses floating-point residuals, which are diagnostics rather than a proof of
MIP optimality. General certificates of branch-and-bound correctness address a much
stronger task \cite{cheung2017certificates}.
